% Informal-meeting version of talk_3dpbl_results: minimal text, big figures.
% Build: tectonic talk_3dpbl_results_informal.tex
\documentclass[11pt,aspectratio=169]{beamer}
\usetheme{default}
\usecolortheme{seahorse}
\setbeamertemplate{navigation symbols}{}
\setbeamertemplate{footline}[frame number]
\setbeamerfont{frametitle}{series=\bfseries,size=\large}
\setbeamertemplate{itemize items}[circle]
\usepackage{amsmath,booktabs}
\graphicspath{{/gpfs/data/fs72996/ewahl/plot_output/diagnostics/}}

\title{Simulating valley turbulence in three dimensions}
\subtitle{First results from the Inn Valley}
\author{Elias Wahl}
\institute{University of Innsbruck}
\date{September 2026}

\begin{document}

\begin{frame}
  \titlepage
\end{frame}

% ---------------------------------------------------------------- scheme
\begin{frame}{The 3D closure in three sentences}
  \begin{itemize}
    \setlength{\itemsep}{2ex}
    \item \textbf{A standard boundary-layer scheme} treats every grid column
      as flat and uniform: it mixes only up--down, and the only wind shear it
      feels is $\partial U/\partial z$. A fine assumption at 10\,km grid
      spacing over plains.
    \item \textbf{In a 500\,m Alpine grid it breaks}: the big eddies are
      half resolved, and on a $30^\circ$ slope a downslope jet shears the
      flow in directions a column scheme cannot see --- most of the
      turbulence production hides there.
    \item \textbf{The 3D closure keeps everything}: same turbulence-energy
      budget as MYNN, but all six stresses and all three components of each
      flux, computed together from \emph{all} the gradients (one small linear
      system per grid point), and mixing applied in all three directions ---
      it replaces the vertical scheme \emph{and} the horizontal diffusion.
  \end{itemize}
\end{frame}

% ---------------------------------------------------------------- changes
\begin{frame}{What we changed since the fork}
  \begin{enumerate}
    \item \textbf{A physically possible answer at every grid point} ---
      no impossible turbulence states, no crashes.
    \item \textbf{Eddy size set by the turbulence itself} --- ended a
      nocturnal runaway.
    \item \textbf{Closed energy bookkeeping} --- slope production is now paid
      by the resolved wind.
    \item \textbf{Host-model bugs separated from scheme physics} --- e.g.\
      radiation cooling shaded ground at 80\,K/h.
    \item \textbf{Exact online budgets} (WRFlux + a TKE budget) --- every
      K/h attributable to a named process; found the energy leak and the
      cold pool's numerical warming.
  \end{enumerate}
\end{frame}

% ---------------------------------------------------------------- result 1
\begin{frame}{The daytime valley wind}
  \begin{columns}[c]
    \column{0.72\textwidth}
    \includegraphics[width=\textwidth]{surface_wind/wind_kolsass_base_vs_3d.png}
    \column{0.26\textwidth}
    Observed maximum $\sim$7\,m/s (10\,m), 10.6 (100\,m).
    \vspace{1ex}

    3D: 6.4 / 10.2.
    \vspace{1ex}

    MYNN: 4.5 / 7--8.
  \end{columns}
\end{frame}

% ---------------------------------------------------------------- result 2
\begin{frame}{One day, every metric (18 July)}
  \begin{center}
  \begin{tabular}{@{}lrrcrr@{}}
    \toprule
    & \multicolumn{2}{c}{lowest 500\,m AGL} & &
      \multicolumn{2}{c}{lowest 1\,km AGL} \\
    \cmidrule{2-3}\cmidrule{5-6}
    & MYNN & 3D & & MYNN & 3D \\
    \midrule
    $\theta$ bias, all launches [K] & 1.63 & \textbf{0.82} & & 1.48 & \textbf{0.86} \\
    $\theta$ bias, day [K] & 2.10 & \textbf{0.29} & & 1.79 & \textbf{0.47} \\
    $\theta$ bias, evening [K] & 0.77 & \textbf{0.17} & & 0.73 & \textbf{0.18} \\
    wind-vector RMSE, soundings [m/s] & 2.20 & \textbf{2.08} & & 2.34 & \textbf{2.21} \\
    wind-speed RMSE, lidar [m/s] & 2.09 & \textbf{1.68} & & 1.89 & \textbf{1.60} \\
    wind-speed bias, lidar [m/s] & $-0.56$ & $\mathbf{-0.02}$ & & $-0.64$ & $\mathbf{-0.17}$ \\
    direction RMSE, lidar [deg] & \textbf{29.7} & 33.2 & & \textbf{28.9} & 30.7 \\
    TKE, $\log_{10}$(model/lidar), day & $-1.15$ & $\mathbf{-0.83}$ & & $-1.29$ & $\mathbf{-0.89}$ \\
    mixed-layer depth error [m] & 207 & \textbf{172} & & 207 & \textbf{172} \\
    \bottomrule
  \end{tabular}
  \end{center}
\end{frame}

% ---------------------------------------------------------------- result 3
\begin{frame}{The model resolves its own convection}
  \begin{columns}[c]
    \column{0.70\textwidth}
    \includegraphics[width=\textwidth]{gz_partition/gz_partition_18.png}
    \column{0.28\textwidth}
    Total turbulence energy equal to MYNN's ---
    \vspace{1ex}

    but 85--91\,\% of it resolved (MYNN: under half).
  \end{columns}
\end{frame}

% ---------------------------------------------------------------- result 4
\begin{frame}{The night-time cold pool: not the scheme's bias}
  \begin{columns}[c]
    \column{0.46\textwidth}
    \includegraphics[width=\textwidth]{x10_fig1_theta_kolsass_05.png}
    \column{0.50\textwidth}
    Pre-dawn pool $\sim$3.5\,K too warm --- \textbf{identically} for both
    schemes, any start time, any forcing resolution.
    \vspace{1ex}

    The measured ingredients: a numerical smoothing term, missing evening
    cold-air transport, a surface layer that will not go calm.
  \end{columns}
\end{frame}

% ---------------------------------------------------------------- outlook
\begin{frame}{Next}
  \begin{itemize}
    \item Twin chains on the driving data (vertical resolution, ICON terrain)
      are running --- the remaining daytime wind deficit is inherited from
      the forcing.
    \item Cold-pool fixes exist per ingredient, all default-off until judged
      against the full observation set.
    \item 20\,m LES of the same transect --- the DKRZ pitch on the next
      slide.
  \end{itemize}
\end{frame}

% ---------------------------------------------------------------- DKRZ
\begin{frame}{The DKRZ pitch: a 20\,m LES of the transect}
  \textbf{The ask: \alert{50\,000 node-hours} on Levante} (+ $\approx$60\,TB
  work / 15\,TB archive). Why we think it flies:
  \vspace{1ex}
  \begin{itemize}
    \item It serves the \textbf{DFG-funded} Tübingen UAS project (Platis):
      virtual drone flights through the LES give the correction functions for
      their measured TKE --- and our grey-zone reference for free.
    \item Nothing speculative: the 500\,m driver is validated and running,
      Gionata's 100\,m LES of the same valley exists, throughput is
      \emph{measured}, not estimated.
    \item The sum: chain 500\,$\to$\,100\,$\to$\,20\,m at $\approx$160
      node-h per simulated hour; six full diurnal cycles (nights included ---
      that's where the deficits live) + sensitivity + margin $=$ 50\,k.
  \end{itemize}
  \vspace{1ex}
  Full arithmetic and the fallback (42\,k):
  \texttt{concept/dkrz\_levante\_antrag\_bullets.md}.
\end{frame}

\end{document}
